\documentclass[reqno]{amsart}
\usepackage[T1]{fontenc}
\usepackage{amsmath,amssymb,amsthm}
\usepackage{xr-hyper}
\usepackage[hidelinks]{hyperref}
\externaldocument[ec-]{20261124T223000--elliptic-curves-for-the-proof__workbook}
\externaldocument[kl-]{20261125T009000--kronecker-limit-and-chowla-selberg__workbook}
\externaldocument[ms-]{20261125T002000--definite-theta-lifts__workbook}
\externaldocument[sw-]{20261125T015500--siegel-weil-in-the-smallest-case__workbook}
\externaldocument[tn-]{20261125T002100--definite-theta-lifts-and-norms__workbook}
\externaldocument[cv-]{20261125T006000--covolumes__workbook}
\externaldocument[cf-]{20261125T017500--chow-forms-and-effective-hilbert-samuel__workbook}
\externaldocument[cp-]{20261125T006500--covolumes-and-periods__workbook}
\newcounter{chap}
% Chapter numbers, generated from 20261125T012000--grh-workbook__guide.org
% by scripts/make-chapters.py.  Do not edit: change the order of the
% index and rerun the script (./preview-site.sh and CI do this).
%
% Usage in a sheet, after \newcounter{chap}:
%   % Chapter numbers, generated from 20261125T012000--grh-workbook__guide.org
% by scripts/make-chapters.py.  Do not edit: change the order of the
% index and rerun the script (./preview-site.sh and CI do this).
%
% Usage in a sheet, after \newcounter{chap}:
%   % Chapter numbers, generated from 20261125T012000--grh-workbook__guide.org
% by scripts/make-chapters.py.  Do not edit: change the order of the
% index and rerun the script (./preview-site.sh and CI do this).
%
% Usage in a sheet, after \newcounter{chap}:
%   \input{chapters}\chapnumber{<sheet ID>}
\makeatletter
\providecommand\chapters@set[2]{\@namedef{chapters@#1}{#2}}
\chapters@set{20261125T001000}{1}% zeros-to-ramp-sums
\chapters@set{20261125T001100}{2}% inert-primes-and-estimate-b
\chapters@set{20261125T005000}{3}% galois-fourier-and-averaging
\chapters@set{20261125T005500}{4}% columns-from-a-sphere
\chapters@set{20261125T000000}{5}% slope-method
\chapters@set{20261126T030100}{6}% a-twist-for-each-prime
\chapters@set{20261125T003000}{7}% the-chain-and-the-square-root
\chapters@set{20261125T010000}{8}% rank-lemma-and-zero-estimate
\chapters@set{20261125T010500}{9}% poincare-series-and-near-orthogonality
\chapters@set{20261124T223000}{10}% elliptic-curves-for-the-proof
\chapters@set{20261125T009500}{11}% fock-model-and-theta-intertwiners
\chapters@set{20261124T223100}{12}% complex-multiplication-and-inert-reduction
\chapters@set{20261125T008000}{13}% one-cm-determinant
\chapters@set{20261124T214500}{14}% supersingular-theta-divisibility
\chapters@set{20261124T214600}{15}% the-isotropic-plane
\chapters@set{20261125T013700}{16}% pairings-on-a-formal-group-of-height-two
\chapters@set{20261125T013500}{17}% hecke-nondegeneracy-on-the-tree
\chapters@set{20261126T030200}{18}% the-bound-is-attained
\chapters@set{20261125T009000}{19}% kronecker-limit-and-chowla-selberg
\chapters@set{20261125T002000}{20}% definite-theta-lifts
\chapters@set{20261125T015500}{21}% siegel-weil-in-the-smallest-case
\chapters@set{20261125T002100}{22}% definite-theta-lifts-and-norms
\chapters@set{20261125T016500}{23}% hecke-operators-and-theta-lifts
\chapters@set{20261125T016700}{24}% divisibility-along-the-chain
\chapters@set{20261125T004000}{25}% denominators-and-sha
\chapters@set{20261125T004100}{26}% tate-twists-and-sha
\chapters@set{20261126T020430}{27}% the-minus-part-by-gauss-sums
\chapters@set{20261125T014100}{28}% congruence-modules
\chapters@set{20261125T014500}{29}% eisenstein-congruences-and-selmer-classes
\chapters@set{20261125T014700}{30}% pseudocharacters-and-the-length-bound
\chapters@set{20261125T011000}{31}% galois-representations-black-box
\chapters@set{20261126T030300}{32}% the-data
\chapters@set{20261126T020210}{33}% bessel-poincare-series-and-their-gram-operators
\chapters@set{20261126T030400}{34}% the-seesaw-identity
\chapters@set{20261126T030500}{35}% divisibility-of-the-composite-lift
\chapters@set{20261126T030600}{36}% theorem-a-short-forms
\chapters@set{20261126T030700}{37}% the-determinant-of-a-chain
\chapters@set{20261126T030800}{38}% the-cost-of-one-lift
\chapters@set{20261126T030900}{39}% denominators-of-projections-to-lifts
\chapters@set{20261126T031000}{40}% twists-with-few-congruences
\chapters@set{20261126T031100}{41}% theorem-c-and-the-main-inequality
\providecommand\chapnumber[1]{%
  \@ifundefined{chapters@#1}%
    {\PackageError{chapters}{Sheet #1 is not listed in the index}%
      {List it in 20261125T012000--grh-workbook__guide.org\MessageBreak
       and run scripts/make-chapters.py.}}%
    {\edef\thechap{\@nameuse{chapters@#1}}}}
\makeatother
\chapnumber{<sheet ID>}
\makeatletter
\providecommand\chapters@set[2]{\@namedef{chapters@#1}{#2}}
\chapters@set{20261125T001000}{1}% zeros-to-ramp-sums
\chapters@set{20261125T001100}{2}% inert-primes-and-estimate-b
\chapters@set{20261125T005000}{3}% galois-fourier-and-averaging
\chapters@set{20261125T005500}{4}% columns-from-a-sphere
\chapters@set{20261125T000000}{5}% slope-method
\chapters@set{20261126T030100}{6}% a-twist-for-each-prime
\chapters@set{20261125T003000}{7}% the-chain-and-the-square-root
\chapters@set{20261125T010000}{8}% rank-lemma-and-zero-estimate
\chapters@set{20261125T010500}{9}% poincare-series-and-near-orthogonality
\chapters@set{20261124T223000}{10}% elliptic-curves-for-the-proof
\chapters@set{20261125T009500}{11}% fock-model-and-theta-intertwiners
\chapters@set{20261124T223100}{12}% complex-multiplication-and-inert-reduction
\chapters@set{20261125T008000}{13}% one-cm-determinant
\chapters@set{20261124T214500}{14}% supersingular-theta-divisibility
\chapters@set{20261124T214600}{15}% the-isotropic-plane
\chapters@set{20261125T013700}{16}% pairings-on-a-formal-group-of-height-two
\chapters@set{20261125T013500}{17}% hecke-nondegeneracy-on-the-tree
\chapters@set{20261126T030200}{18}% the-bound-is-attained
\chapters@set{20261125T009000}{19}% kronecker-limit-and-chowla-selberg
\chapters@set{20261125T002000}{20}% definite-theta-lifts
\chapters@set{20261125T015500}{21}% siegel-weil-in-the-smallest-case
\chapters@set{20261125T002100}{22}% definite-theta-lifts-and-norms
\chapters@set{20261125T016500}{23}% hecke-operators-and-theta-lifts
\chapters@set{20261125T016700}{24}% divisibility-along-the-chain
\chapters@set{20261125T004000}{25}% denominators-and-sha
\chapters@set{20261125T004100}{26}% tate-twists-and-sha
\chapters@set{20261126T020430}{27}% the-minus-part-by-gauss-sums
\chapters@set{20261125T014100}{28}% congruence-modules
\chapters@set{20261125T014500}{29}% eisenstein-congruences-and-selmer-classes
\chapters@set{20261125T014700}{30}% pseudocharacters-and-the-length-bound
\chapters@set{20261125T011000}{31}% galois-representations-black-box
\chapters@set{20261126T030300}{32}% the-data
\chapters@set{20261126T020210}{33}% bessel-poincare-series-and-their-gram-operators
\chapters@set{20261126T030400}{34}% the-seesaw-identity
\chapters@set{20261126T030500}{35}% divisibility-of-the-composite-lift
\chapters@set{20261126T030600}{36}% theorem-a-short-forms
\chapters@set{20261126T030700}{37}% the-determinant-of-a-chain
\chapters@set{20261126T030800}{38}% the-cost-of-one-lift
\chapters@set{20261126T030900}{39}% denominators-of-projections-to-lifts
\chapters@set{20261126T031000}{40}% twists-with-few-congruences
\chapters@set{20261126T031100}{41}% theorem-c-and-the-main-inequality
\providecommand\chapnumber[1]{%
  \@ifundefined{chapters@#1}%
    {\PackageError{chapters}{Sheet #1 is not listed in the index}%
      {List it in 20261125T012000--grh-workbook__guide.org\MessageBreak
       and run scripts/make-chapters.py.}}%
    {\edef\thechap{\@nameuse{chapters@#1}}}}
\makeatother
\chapnumber{<sheet ID>}
\makeatletter
\providecommand\chapters@set[2]{\@namedef{chapters@#1}{#2}}
\chapters@set{20261125T001000}{1}% zeros-to-ramp-sums
\chapters@set{20261125T001100}{2}% inert-primes-and-estimate-b
\chapters@set{20261125T005000}{3}% galois-fourier-and-averaging
\chapters@set{20261125T005500}{4}% columns-from-a-sphere
\chapters@set{20261125T000000}{5}% slope-method
\chapters@set{20261126T030100}{6}% a-twist-for-each-prime
\chapters@set{20261125T003000}{7}% the-chain-and-the-square-root
\chapters@set{20261125T010000}{8}% rank-lemma-and-zero-estimate
\chapters@set{20261125T010500}{9}% poincare-series-and-near-orthogonality
\chapters@set{20261124T223000}{10}% elliptic-curves-for-the-proof
\chapters@set{20261125T009500}{11}% fock-model-and-theta-intertwiners
\chapters@set{20261124T223100}{12}% complex-multiplication-and-inert-reduction
\chapters@set{20261125T008000}{13}% one-cm-determinant
\chapters@set{20261124T214500}{14}% supersingular-theta-divisibility
\chapters@set{20261124T214600}{15}% the-isotropic-plane
\chapters@set{20261125T013700}{16}% pairings-on-a-formal-group-of-height-two
\chapters@set{20261125T013500}{17}% hecke-nondegeneracy-on-the-tree
\chapters@set{20261126T030200}{18}% the-bound-is-attained
\chapters@set{20261125T009000}{19}% kronecker-limit-and-chowla-selberg
\chapters@set{20261125T002000}{20}% definite-theta-lifts
\chapters@set{20261125T015500}{21}% siegel-weil-in-the-smallest-case
\chapters@set{20261125T002100}{22}% definite-theta-lifts-and-norms
\chapters@set{20261125T016500}{23}% hecke-operators-and-theta-lifts
\chapters@set{20261125T016700}{24}% divisibility-along-the-chain
\chapters@set{20261125T004000}{25}% denominators-and-sha
\chapters@set{20261125T004100}{26}% tate-twists-and-sha
\chapters@set{20261126T020430}{27}% the-minus-part-by-gauss-sums
\chapters@set{20261125T014100}{28}% congruence-modules
\chapters@set{20261125T014500}{29}% eisenstein-congruences-and-selmer-classes
\chapters@set{20261125T014700}{30}% pseudocharacters-and-the-length-bound
\chapters@set{20261125T011000}{31}% galois-representations-black-box
\chapters@set{20261126T030300}{32}% the-data
\chapters@set{20261126T020210}{33}% bessel-poincare-series-and-their-gram-operators
\chapters@set{20261126T030400}{34}% the-seesaw-identity
\chapters@set{20261126T030500}{35}% divisibility-of-the-composite-lift
\chapters@set{20261126T030600}{36}% theorem-a-short-forms
\chapters@set{20261126T030700}{37}% the-determinant-of-a-chain
\chapters@set{20261126T030800}{38}% the-cost-of-one-lift
\chapters@set{20261126T030900}{39}% denominators-of-projections-to-lifts
\chapters@set{20261126T031000}{40}% twists-with-few-congruences
\chapters@set{20261126T031100}{41}% theorem-c-and-the-main-inequality
\providecommand\chapnumber[1]{%
  \@ifundefined{chapters@#1}%
    {\PackageError{chapters}{Sheet #1 is not listed in the index}%
      {List it in 20261125T012000--grh-workbook__guide.org\MessageBreak
       and run scripts/make-chapters.py.}}%
    {\edef\thechap{\@nameuse{chapters@#1}}}}
\makeatother

\makeatletter
\@ifundefined{chapters@20261125T020300}{\renewcommand{\thechap}{?}}{\chapnumber{20261125T020300}}
\makeatother
\renewcommand{\thesection}{\thechap.\arabic{section}}
\theoremstyle{definition}
\newtheorem{problem}{Problem}[section]
\newtheorem{blackbox}{Input}[section]
\newcommand{\rd}{\operatorname{rd}}
\newcommand{\mass}{\operatorname{mass}}
\newcommand{\htop}{\operatorname{ht}}
\newcommand{\lcv}{\operatorname{lcv}}
\newcommand{\Q}{\mathbb{Q}}
\setlength{\emergencystretch}{2em}
\title{Heights, the CM period and the auxiliary prime}
\author{}
\date{November 25, 2026}
\begin{document}
\maketitle
\refstepcounter{chap}\label{chap}

\paragraph{Prerequisites.}
The period calculation in Chapters~\ref{ec-chap} and~\ref{kl-chap},
the mass and the squared norm of a lift in Chapters~\ref{ms-chap}
and~\ref{tn-chap}, and the degrees of lattices of sections in
Chapters~\ref{cv-chap}, \ref{cf-chap} and~\ref{cp-chap}.
We use squared operator norms and slopes equal to twice the arithmetic
degree divided by the field degree and the rank \cite[\S15.1]{Adr}.

\section{The period in a squared norm}\label{sec:period}

Which linear term accompanies $a\log a-a$ in the height of a theta lift?
For a complex lattice $L$ write $\mathcal A(L)=\operatorname{area}(L)
|\Delta(L)|^{1/6}$, where $\Delta=g_2^3-27g_3^2$.
This quantity is invariant under homothety.  Let $d\equiv7\pmod{16}$
be prime, $K=\Q(\sqrt{-d})$, and $\chi_d$ the primitive character of
discriminant $-d$.  Let $E\supset K$ be a CM field whose relative different
is $\mathfrak b\bar{\mathfrak b}$, and let $\Phi$ consist of the
$f=[E:K]$ embeddings extending a fixed complex embedding of $K$.
Write $v=\rd_E/2$, where $\rd_E=|\operatorname{disc}E|^{1/(2f)}$,
and $c_{\chi_d}=L'(1,\chi_d)/L(1,\chi_d)$.

\begin{blackbox}[CM period]\label{in:period}
Let $L_C$ represent each of the $h_K$ ideal classes of $K$ as a complex lattice.
At each $\tau\in\Phi$, the normalized arithmetic degree $\log\Omega_\tau$
of the coordinate line of periods $\mathcal L_\tau$ over a field of definition
$F$ means $2\widehat{\deg}(\mathcal L_\tau)/[F:\Q]$ \cite[Def.~15.5]{Adr}.  Its value is
\[
 \log\Omega_\tau
 =\frac1{h_K}\sum_C\log\frac{\mathcal A(L_C)}{\pi}
       -\log\frac{\rd_E}{\sqrt d}
 =\gamma+\log\frac{\pi}{v}-c_{\chi_d} =:\log\Omega .
\]
Simplest incarnation: Problem~\ref{kl-pr:LL-eta-gen} computes the class
mean by Kronecker's limit formula (Theorem~\ref{kl-thm:klf}).
\emph{Status}: the class mean is classical, by Kronecker's limit formula;
the relative-different correction is proved in \cite[Props.~15.3 and~15.6]{Adr}.
\end{blackbox}

\begin{problem}\label{pr:quadratic}
For $E=K$ the correction is zero and
$\log\Omega=\gamma+\log(2\pi/\sqrt d)-c_{\chi_d}$.
Using $\mathcal A(\mathbb Z+\mathbb Z\tau)=4\pi^2\operatorname{Im}(\tau)
|\eta(\tau)|^4$ and the Gamma product of Theorem~\ref{kl-thm:cs} gives
$\log\Omega=-\tfrac12\log d+
\tfrac{w_K}{2h_K}\sum_{j=1}^{d-1}\chi_d(j)\log\Gamma(j/d)$,
where $w_K$ counts the roots of unity in $K$.
\end{problem}

\begin{problem}\label{pr:gaussian}
In the Gaussian example, $\Omega_{\mathrm i}=4\pi|\eta(\mathrm i)|^4
=\Gamma(1/4)^4/(4\pi^2)$ (Problem~\ref{ec-pr:Omega}).
Thus $\log\Omega_{\mathrm i}=1.4763\ldots$, whereas the bound for a single
Gaussian lift in Theorem~\ref{tn-thm:norm} is
$\log\|\Theta_{(a)}\|^2\le a\log a-a-a\log(2\pi)+O(\log a)$.
Multiplying its squared norm by $(2\pi/\Omega_{\mathrm i})^a$
changes the last linear term to $-a\log\Omega_{\mathrm i}$.
\end{problem}

\section{The height and the target slope}\label{sec:height}

Fix one translate and its nodes throughout.  For large $x$ and fixed
$0<\varepsilon<1/52$ write $y=x^\varepsilon$, $b=e^y$, $a=xb$,
$\Lambda_a=\log(a\rd_E)$ and $\varphi(z)=z\log z-z$, with $\varphi(0)=0$.
The chain has $m$ lifts, initial rank $k$, node weights $w_{j\rho}>0$
($1\le j\le m$, $\rho\in\Phi$), and $N_T$ distinguished embeddings
(the tops $t\in T$) with weights $a_t$ and sum $A=\sum_t a_t$.
A word $u$ specifies a level prime $T_u$, a source lattice $N_u$ and a
period map $B_u=P_u K_m\cdots K_1:N_u\to R^\vee$, where $P_u$ is
the final pairing and $R$ is the top lattice.  Over the field of definition,
$\htop(B_u)$ is the mean of $\log\|B_u\|^2$ over the archimedean places.

\begin{blackbox}[Height of the composite]\label{in:height}
Let $s_j$ bound the mean log squared norm of $K_j$.  In the frames of the
chain the specified bound is
\[
 s_j=\sum_{\rho,i}\bigl(\varphi(\nu_{j\rho i})-
                  \nu_{j\rho i}\log\Omega\bigr)
       + C_0fN^2(\Lambda_a+m^3y),
\]
where $C_0\ge1$ is independent of the word and $\nu_{j\rho i}$ are the $N$ row weights at $\rho$.
Write $h_{j\rho}$ for the log hook constant of its frame,
$e_{\lambda,j}^{\max}$ for its error allowance, $e_P^{\max}$ for
the allowance in $\log\|P_u\|^2\le-A\log T_u+e_P^{\max}$,
and $\langle\log\kappa\rangle$ for the mean log scalar of the target
metric.  Use the allowances of \cite[Def.~15.14]{Adr}.  Under the kernel, pairing and frame bounds of
\cite[Hyps.~15.10--15.12]{Adr}, for $B_u\ne0$,
\[
 \htop(B_u)\le h_u:=-A\log T_u+e_P^{\max}-\langle\log\kappa\rangle
       +\sum_j\bigl(s_j-\sum_\rho h_{j\rho}+e_{\lambda,j}^{\max}\bigr).
\]
Simplest incarnation: the squared norm in Theorem~\ref{tn-thm:norm}
has one factorial term.  \emph{Status}: composition in the stated frames is proved in
\cite[Prop.~15.9, Cor.~15.15]{Adr}.  The bounds for the chain are proved in
the Lean formalization \cite{Lean}; they are \cite[Hyps.~15.10--15.12]{Adr},
not yet written out in the published expositions.
\end{blackbox}

\begin{problem}\label{pr:common-height}
Suppose the terms defining $h_u$ other than $-A\log T_u$ are common
to all words, as in \cite[Prop.~18.11(vi)]{Adr}.
Then $(h_u+A\log T_u)-(h_{u'}+A\log T_{u'})=0$.
Condition (K6) only bounds the absolute value of this difference.
\end{problem}

The target slope is bounded through the covolumes of lattices $M_j$ of
forms on unitary groups of signatures $(p_j,q_j)=(k+j-1,m-j)$.
Let $D=\operatorname{rk} R$, $n=f/2$, $N_2$ the size of $T_2$, the set indexing the thick weights, $B_0=fmby$ and $N_c=N_2kby^2/n^2$.
Write $\mathcal L$ for the sum of the normalized local lengths at the
twisted places, as in \cite[\S15.4]{Adr}.

\begin{blackbox}[Covolumes and the height inequality]\label{in:A4}
The covolume estimate on the chain, with $\lcv=-2\widehat{\deg}/[F:\Q]$ over the field of definition $F$, is
\[
 \frac{\lcv(M_j)}D\le
 \sum_{\rho}(p_j-q_j)w_{j\rho}\log\Omega+C B_0\log\log x .
\]
For sufficiently small $\varepsilon$ and each admissible family there are
$C,x_0$ for which this holds at every $x\ge x_0$, uniformly in the
point, lift and level.  After passage between levels and the local
comparisons, the target bound of \cite[Hyp.~15.16]{Adr} reads
\[
 \mu(R^\vee)-\langle\log\kappa\rangle
 \le \sum_{j,\rho}(p_j-q_j)w_{j\rho}\log\Omega+\mathcal L
       +1+C_0 b+C_0my(B_0+N_c).
\]
With Input~\ref{in:height}, the kernel weight relations and the size
conditions of \cite[Hyp.~15.18, Thm.~15.19]{Adr}, this gives (A4):
\begin{align*}
 \mu(R^\vee)+h_u-\mathcal L
   &\le m\sum_t\bigl(\varphi(a_t)-a_t\log\Omega\bigr)-A\log T_u+e_4,\\
 e_4&=(50C_0+10)\bigl(N_Ta\Lambda_a+my(B_0+N_c)\bigr).
\end{align*}
Simplest incarnation: for theta sections on an elliptic curve $J$ with
a unit-discriminant model, Theorem~\ref{cp-thm:Mk} gives the slope
$-h_F(J)+\tfrac12\log(k/(2\pi^2))$ in weight $k$, with $h_F$
the Faltings height.  \emph{Status}: proved in the Lean formalization; for
the chain it rests on \cite[Hyps.~15.10--15.12,
15.16, 15.18, 15.58--15.59, 15.69]{Adr}, not yet written out in the published expositions.
\end{blackbox}

\begin{problem}\label{pr:weights}
At $(j,\rho)$ the kernel weight has sum
$\mathbf1_T(\rho)a_\rho+kw_{j\rho}+\delta_{j\rho}$, with the top
term taken as zero off $T$.  After adding the covolume term, the coefficient
of $w_{j\rho}\log\Omega$ is $2j-1-m$, of absolute value at most $m-1$.
The top terms over all lifts contribute $-mA\log\Omega$.
\end{problem}

\begin{problem}\label{pr:error}
If $\Lambda_a\le2y$ as in \cite[Lemma~18.15(iii)]{Adr} and $y\ge1$, then
$e_4\le(100C_0+20)y\bigl(N_Ta+m(B_0+N_c)\bigr)$,
which also gives the bound with $y^2$ used in the assembly.
\end{problem}

\begin{problem}\label{pr:uniform}
On the fixed projective line, Theorem~\ref{cv-thm:P1} gives
$\deg H^0(\mathcal O(D))=D^2/4+O(D\log D)$.
For the two points cut out by $F(X,Y)=X(X-HY)$, $H$ a positive integer,
Theorem~\ref{cf-thm:points} gives
$\deg\Lambda_D=\tfrac D2\log(1+H^2)-\log H$ for $D\ge1$,
where $\Lambda_D$ is the lattice of sections on these points with
the evaluation metric.  The constant term is unbounded in absolute
value as $H$ varies.
\end{problem}

For the chain, \cite[Thms.~15.46 and~15.51]{Adr} reduce the norm estimates to
the normalized kernel, theta-null mean and class counts
\cite[Hyps.~15.42--15.44]{Adr}, and to the branching norm
\cite[Hyp.~15.50]{Adr} together with factorization of the final pairing, its evaluation bound
and the stated weight and size conditions.
Siegel--Weil expresses a genus average by a product of local densities
in Theorem~\ref{sw-thm:rank8-16}.  With the doubling method it gives the
Rallis inner product formula in the convergent range.  For the frame
constants one still needs the Petersson normalization and a lower bound
for its factors \cite[Prop.~15.52, \S18.7]{Adr}.
For the target slope, \cite[Thm.~15.70]{Adr} uses passage, the base-side
comparison and local estimates.  The passage uses near-constant
multiplicities \cite[Hyp.~15.56]{Adr}.  The covolume deduction
\cite[Thm.~15.66]{Adr} assumes (I1)--(I17) and a uniform mean slope
on deep covers \cite[Hyp.~15.69]{Adr}.  Arithmetic Hilbert--Samuel
on a fixed model (Theorem~\ref{cv-thm:HS}) gives asymptotics there.
The component estimates needed to pass to a mean slope are spelled out in
\cite[Lemma~15.62]{Adr}.  The kernel weight relations remain a separate
hypothesis \cite[Hyp.~15.18]{Adr}.

\section{Counting forms and choosing the prime}\label{sec:prime}

How large a prime can the final pairing use?  Let $m\ge2$ be even,
$r=k+m-1$ and $g=\log(g_0g_1)$ for the two level primes $g_0,g_1$.
For a shifted highest weight $z$ of length $s$ let
$V_s(z)=\prod_{i<i'}(z_i-z_{i'})$ and $D_s(z)=V_s(z)/\prod_{i=0}^{s-1}i!$.
For the genus of the standard rank-$s$ hermitian lattice write $\mass(s)$
for the sum of reciprocal stabilizer orders.  Let $I_\Pi,I_\Sigma$ be the
source and target level indices at $g_0,g_1$.  For the column and row
parameters $z^{\rm col}_{u\rho},z^{\rm row}_\rho$ define
\[
 d_u^*=\mass(k) I_\Pi\prod_{\rho}D_k(z^{\rm col}_{u\rho}),
 \qquad X_u=(D^*/d_u^*)^{1/(fk)},
\]
where $D^*$ is the target dimension budget.  The exponent $fk$ comes from
the source dimension $d_u^*(T_u^k-1)^f$ for trivial stabilizers.

\begin{blackbox}[The dimension budget and the level]\label{in:A6}
Assume the row and column interlacing, mass-ratio formula and size
conditions of \cite[Prop.~15.25]{Adr}, and the dimension bound
\[
 D^*\ge(2/3)^f\mass(r)I_\Sigma\prod_\rho D_r(z^{\rm row}_\rho).
\]
Let $\Gamma_0\ge1$ be the frame factor, $\log\Gamma_0\le fk$, and let
$T_u\ge Z_u:=X_u\Gamma_0^{-2/(fk)}$.  Then (A6) is
\[
 \log T_u\ge\sum_{j=1}^m\frac1f\sum_{\rho\in\Phi}
                    \log\frac{w_{j\rho}v}{\pi}+mg-e_6,
 \qquad e_6=C_6y(1+m^2/k)+2 .
\]
Simplest incarnation: Theorem~\ref{ms-thm:count} cancels the Weyl
factorials in the Gaussian mass, leaving $V_s(z)\pi^{-s(s+1)/2}$
times the local factors.  \emph{Status}: the deduction is proved in
\cite[Prop.~15.25, Thm.~15.26]{Adr}; the target budget is proved in the
Lean formalization, and is \cite[Hyp.~15.23]{Adr}, not yet written out in the published expositions.  The mass and formal degree must
use the same measures \cite[\S18.7]{Adr}.
\end{blackbox}

\begin{problem}\label{pr:factorials}
For each $k<M\le r$ the mass ratio contributes $((M-1)!)^f$.
The Weyl denominators in the row-to-column ratio cancel their product
exactly.  A row with $k-1$ entries $v_i\ge0$ and $m$ entries
$-b_j<0$ contributes at least $b_j^{k-1}$ from the cross differences
at node $j$.  After division of logarithms by $fk$ the coefficient of
the mean $\log b_j$ from these factors is $(k-1)/k$.
\end{problem}

\begin{problem}\label{pr:indices}
At a level prime $g_i$ the leading powers in the target and source
indices are $g_i^{r(r+1)/2+k}$ and $g_i^{k(k+1)/2}$ respectively.
Each index lies between one half and one times its leading power
\cite[Lemma~15.21]{Adr}.  Taking both primes at each of the $f$
embeddings gives
\[
 \frac1{fk}\log\frac{I_\Sigma}{I_\Pi}
 \ge \left(m+\frac{m(m-1)}{2k}\right)g-\frac{2\log2}{k}.
\]
\end{problem}

\begin{problem}\label{pr:interval}
For $Z_u\le T_u\le(1+1/(fk))Z_u$,
\[
 D^*\Gamma_0^{-2}\le d_u^* T_u^{fk}\le e D^*\Gamma_0^{-2},
 \qquad \log T_u\ge\log X_u-2 .
\]
\end{problem}

The required congruence is $T_u=\alpha\bar\alpha$ in $K$, with
$\alpha\equiv1\pmod {M\mathcal O_K}$ for
$M=Q_Sqd(g_0g_1)^2$ \cite[Rem.~15.27]{Adr}.  Here $q$ is the first
auxiliary character modulus and $Q_S$ is the part of the conductor
of $E$ prime to $qd$.  The ray class estimate in \cite[\S5.8, (C9)]{Adr}, from \cite[Thm.~1.4]{TZ} with Stark's exceptional-zero bound,
supplies such a prime in $(Z,(1+h)Z]$ when
$\log(4dM^2)\le c\sqrt{\log Z}$ and
$e^{-c\sqrt{\log Z}}\le h\le1$, for an absolute $c>0$.

\begin{problem}\label{pr:ray}
Take $m\asymp x^{1/2-\varepsilon/16}$, $fk=x^{O(1)}$,
$\log(Q_Sqd)\ll x^{\varepsilon/5}\log^2 x$, and
$(kb)^8\le g_i\le2(kb)^8$ as in \cite[Prop.~5.15, Def.~6.1]{Adr}.
If $\log Z_u\gg my$, then $\log(4dM^2)\asymp y=o(\sqrt{\log Z_u})$
and $\log(fk)=o(\sqrt{\log Z_u})$.  Hence (C9) supplies $T_u$ in
the interval of Problem~\ref{pr:interval}.
\end{problem}

Under the counting hypotheses the required lower scale follows already
from \cite[Prop.~15.25]{Adr}: $w_{j\rho}\ge1$, $v\ge\pi$,
$g\ge16y$ and $k\ge m$ give
$\log Z_u\ge mg-O(y(1+m^2/k))-2$.  In the ranks of
\cite[Prop.~5.15]{Adr}, $m/k\to0$, so the error is $o(my)$.

\begin{problem}\label{pr:linear}
For each top $t$ and node $j$, subtracting the counted node term from
the height term gives
\[
 \varphi(a_t)-a_t\log\Omega-
 \frac{a_t}{f}\sum_\rho\log\frac{w_{j\rho}v}{\pi}
 =a_t\left(\log a_t-1-\gamma+c_{\chi_d}
               -\frac1f\sum_\rho\log w_{j\rho}\right).
\]
\end{problem}

In the mean slope inequality the average of $\mu(R^\vee)+h_u$ is
bounded using Input~\ref{in:A4} in the same metric.  Multiplying
Input~\ref{in:A6} by $-A$ then contributes $-mAg$, which cancels the
$+mAg$ in the local bound (A5) \cite[Prop.~3.1]{Adr}.
Problem~\ref{pr:linear} gives the coefficient that survives.

\section{The base lattice}\label{sec:base}

Let $M_u$ be the level-zero lattice for the word $u$.  Its coefficient
module at $\rho$ is the integral Weyl module of highest weight
$\lambda_\rho$, with the Fock form normalized to give its highest
weight vector squared norm one.  For a dominant weight $\lambda$ of
length $k$, let $|\lambda|=\sum_i\lambda_i$ and subtract the last
entry from every entry to obtain the partition $\nu$.
Let $\mu_{\rm std}(\nu)$ be the slope of this normalized Weyl lattice.
In the unnormalized Fock form the monomials are orthogonal with
$\|x^\alpha\|^2=\alpha!$.  Let $c_\nu=\|hw_\nu\|^2/\prod_i\nu_i!$.

\begin{blackbox}[Slope of the base lattice]\label{in:A7}
In the level-zero genus,
\[
 \mu(M_u)=-\frac{\log T_u}{k}\sum_\rho |\lambda_\rho|
             +\sum_\rho\mu_{\rm std}(\nu_\rho),\qquad
 \mu_{\rm std}(\nu)\ge-|\nu|\log k-\log c_\nu .
\]
Under the weight and size bounds of \cite[Thm.~15.33]{Adr}, including
$\log T_u\le C_Tmy$, (A7) is $\mu(M_u)\ge-e_7$, with
$e_7=C_7 N_T a\Lambda_a(1+m/k)$.
Simplest incarnation: Problem~\ref{cv-pr:def-n1} computes the slope
of a Gaussian invariant line in weight divisible by four.  Its
Petersson squared norm includes a stabilizer factor $1/4$, giving
slope $\log4$ where the normalized coefficient line has slope zero.  \emph{Status}: the displayed formulas and the lower
bound are proved in \cite[Props.~15.30, 15.32, Thm.~15.33]{Adr}.
\end{blackbox}

\begin{problem}\label{pr:hook}
For $\nu=(l,h)$ with $l\ge h\ge0$, the hook formula of
\cite[Lemma~15.29]{Adr} gives
$\|hw_\nu\|^2=(l+1)!h!/(l-h+1)$ and $c_\nu=(l+1)/(l-h+1)$.
For $(l,h)=(2,1)$ the polynomial
$x_{11}(x_{11}x_{22}-x_{12}x_{21})$ has squared Fock norm $3$.
\end{problem}

\begin{problem}\label{pr:record}
With $y\le\Lambda_a$, Input~\ref{in:A6} gives the lower inequality in
(K7) with error $(C_6+2)\Lambda_a(1+m^2/k)$.
Input~\ref{in:A7} gives the lower inequality in (K8).
The upper bound on $\mu_{\max}(M_u)$ is a separate input to
\cite[Prop.~18.11(viii), Rem.~15.34]{Adr}.
\end{problem}

\paragraph{Deliverable.}
For the chain record of \cite[Def.~18.9]{Adr}, Input~\ref{in:height}
specifies $h_u$, Input~\ref{in:A4} supplies the height inequality in
(K4), and Problem~\ref{pr:common-height} supplies (K6) under the
common-frame assumption.  Inputs~\ref{in:A6} and~\ref{in:A7} supply
the first inequalities of (K7) and (K8).  The upper and between-word
bounds on $\log T_u$ and the maximal base slope still enter
\cite[Prop.~18.11(vii)--(viii)]{Adr} separately.  These are the
contributions to (A4), (A6) and (A7) in \cite[Prop.~18.14]{Adr}.
Existence of records uniformly on the window is proved in the Lean
formalization; it is \cite[Hyp.~18.10]{Adr}, not yet written out in the published expositions.

\section*{Hints}
\subsection*{\S\ref{sec:period}}
\begin{description}
\item[\ref{pr:quadratic}] Take logarithms of the product over classes.
\end{description}
\subsection*{\S\ref{sec:prime}}
\begin{description}
\item[\ref{pr:interval}] Use $\log(1+t)\le t$.
\item[\ref{pr:ray}] Compare $y^2$ with $my$.
\end{description}

\begin{thebibliography}{Adr}
\bibitem[Adr]{Adr}
J.~Adriazola, \emph{The Alp\"oge proof of the Riemann hypothesis}, manuscript, revised 24 November 2026, \url{https://github.com/jimjamzola/RH-Alpoge-Reconstruction}.
\bibitem[Lean]{Lean} The Lean community, \emph{riemann-hypothesis}, a Lean formalization of Alp\"oge's argument, \url{https://github.com/leanprover-community/riemann-hypothesis}.
\bibitem[TZ]{TZ} J.~Thorner and A.~Zaman, \emph{A unified and improved Chebotarev density theorem},
Algebra Number Theory \textbf{13} (2019), 1039--1068.
\end{thebibliography}
\end{document}
